% Build: tectonic PlanBee_Architecture.tex
% Korean font: Apple SD Gothic Neo (macOS); replace with an installed Korean font elsewhere.
\documentclass[11pt,a4paper]{article}
\usepackage[margin=23mm]{geometry}
\usepackage{fontspec}
\setmainfont{Apple SD Gothic Neo}
\setsansfont{Apple SD Gothic Neo}
\setmonofont{Menlo}[Scale=0.85]
\usepackage{amsmath,booktabs,tabularx,array,xcolor,tikz,enumitem,fancyhdr,hyperref}
\usetikzlibrary{arrows.meta,positioning,calc}
\definecolor{ink}{HTML}{18334B}
\definecolor{blue}{HTML}{EAF3FC}
\definecolor{green}{HTML}{EAF5EF}
\definecolor{amber}{HTML}{FFF3DE}
\hypersetup{colorlinks=true,linkcolor=ink,urlcolor=ink}
\XeTeXlinebreaklocale "ko"
\XeTeXlinebreakskip=0pt
\setlength{\parindent}{0pt}
\setlength{\parskip}{6pt}
\setlength{\emergencystretch}{2em}
\setlist{nosep,leftmargin=1.5em}
\pagestyle{fancy}\fancyhf{}\fancyhead[L]{\small Plan Bee / V2 Architecture}\fancyhead[R]{\small 2026-11-23}\fancyfoot[C]{\thepage}
\setlength{\headheight}{15pt}
\newcolumntype{Y}{>{\raggedright\arraybackslash}X}
\tikzset{box/.style={draw=ink!60,rounded corners=3pt,fill=blue,align=center,text width=11.5cm,minimum height=9mm,font=\small,inner sep=5pt},arr/.style={-{Latex[length=2mm]},thick,draw=ink}}

\definecolor{gray}{HTML}{F1F2F4}
\tikzset{
 flow/.style={draw=ink!50,rounded corners=3pt,fill=blue,text width=9.4cm,align=left,inner sep=5pt,font=\small},
 side/.style={draw=ink!40,dashed,rounded corners=3pt,fill=gray,text width=4.7cm,align=left,inner sep=5pt,font=\footnotesize},
 branch/.style={arr,dashed,draw=ink!70}}
% #1 제목(단계 | 주체: 하는 일), #2/#3 첫째 줄 라벨/내용, #4/#5 둘째 줄 라벨/내용
\newcommand{\flowstep}[5]{{\raggedright\textbf{#1}\par\vspace{1pt}{\footnotesize\textbf{#2}\enspace #3\par\textbf{#4}\enspace #5\par}}}
\newcommand{\flowtext}[3]{{\raggedright\textbf{#1}\par\vspace{2pt}{\footnotesize\textbf{입력·판단} #2\par\textbf{결과} #3\par}}}
\newcommand{\sidetext}[2]{{\raggedright\textbf{#1}\par\vspace{1pt}#2\par}}

\begin{document}
\begin{center}
{\Huge\bfseries Plan Bee}\\[5mm]
{\LARGE 전체 실행 구조와 에이전트 협업 흐름}\\[3mm]
{\large 후보 생성 · 위임카드 기반 토론 · PM 최종 결정}\\[3mm]
2026년 11월 23일 / 현행 V2 구현 기준
\end{center}
\vspace{3mm}
\textbf{핵심 구조.} 사람이 업무 변경 이슈를 줄글로 입력하면, 생성 AI와 알고리즘이 새로운 milestone 계획 후보를 만든다. Hive는 일정과 하드 제약을 코드로 검증한다. 각 개인의 Proxy LLM은 자신의 위임카드로 후보들을 평가하고 공개 토론한다. 추천 후보 1--3개를 받은 PM이 최종 선택하며, 필요한 본인 확인과 최신 Hive 검증이 끝나야 실제 보드에 반영된다.

\textbf{이 문서의 범위.} 현재 서비스의 \texttt{candidate\_discussion} 경로를 설명한다. 이전 \texttt{legacy} 협상 경로는 호환용으로 남아 있으나 본문의 주 흐름과 구분한다. 구현 설명과 실제 Live 성능 검증을 구분한다. T1--T5 Live 결과는 저장된 실험 기록(v006--v009)에서 옮겼고, 남은 한계는 마지막 절에 기록한다.

\section{주체와 판단 권한}
\begin{tabularx}{\linewidth}{@{}p{27mm}YY@{}}\toprule
주체 & 수행하는 일 & 수행하지 않는 일\\\midrule
사람 사용자 & 역할 선택, 이슈 입력, 본인 확인, PM 역할의 최종 선택 & 채팅 입력만으로 계획 자동 승인\\
생성 AI & Task 내용에 근거한 분해·이동·지원 구조 설계 & 개인 비공개 카드 열람, 공정성 승자 결정\\
탐색 알고리즘 & 담당·시작·분량 조합 확장, 중복 제거 & 최소 부담 분산만으로 적용안 선정\\
Hive / 엔진 & 일정 계산, 제약 검사, 호출 순서·상태·예산 관리 & 개인의 카드상 의사를 대신 결정\\
개인 Proxy LLM & 자기 카드에 따른 평가·발언·거절·동의 & 다른 사람 카드 원문 열람, 마감 제약 무시\\
PM & 추천 후보 선택·전부 반려·수정 요청 & 당사자 거절 자동 무효화, 하드 제약 우회\\\bottomrule
\end{tabularx}

Hive의 일정 검증은 결정적 코드이며 별도의 LLM 심판이 아니다. 자연어 위임 의사는 LLM이 판단하고, 숫자·일정·구조의 유효성은 코드가 검증한다.

\newpage
\section{전체 실행 흐름}
\textbf{읽는 법.} 왼쪽 열은 정상 경로이며 01부터 09까지 순서대로 진행한다. 상자 색은 처리 주체를 뜻한다. \colorbox{blue}{파랑}은 LLM 또는 탐색 알고리즘, \colorbox{green}{초록}은 결정적 코드(Hive), \colorbox{amber}{노랑}은 사람이다. 오른쪽 \colorbox{gray}{회색 점선 상자}는 조건이 충족될 때만 들어가는 예외 경로다. 09를 통과하기 전까지 모든 후보는 복제 보드 위의 미리보기이며 실제 보드는 바뀌지 않는다.
\begin{center}
\begin{tikzpicture}[node distance=2mm,every node/.append style={inner sep=4pt}]
\node[flow] (input) {\flowstep{01 | 사용자: 이슈 등록, run 시작}{입력}{역할, 대상 업무, 줄글 이슈 1건.}{출력}{새 run. 현재 보드·위임카드·프롬프트 버전을 snapshot으로 고정하며, 이후 단계는 모두 이 snapshot 기준.}};
\node[flow,below=of input] (estimate) {\flowstep{02 | 대상 업무 담당자의 Proxy LLM: 이슈를 일정 변화로 변환}{입력}{담당자 본인 카드, 대상 업무 정보, 이슈 원문.}{출력}{``시작일만 밀림'' 또는 ``소요 시간 증가''와 새 소요일. Hive가 이를 반영해 마감 초과 여부를 계산하고, 초과면 기한 회복 run, 충족이면 공정성 재분담 run으로 정한다.}};
\node[flow,below=of estimate] (generate) {\flowstep{03 | 생성 AI + 탐색 알고리즘: 수정 계획 후보 생성}{입력}{이슈, Task 설명·완료 조건·분리 가능 경계, 공개 일정. 개인 위임카드는 주지 않는다.}{출력}{후보 계획 목록(조합 최대 500개). 각 후보는 업무 분해·담당 변경·시작일 이동·지원 배정·Beta 이관을 조합한 새 보드다.}};
\node[flow,fill=green,below=of generate] (validate) {\flowstep{04 | Hive(코드): 후보마다 제약 검사}{검사}{후보를 적용한 복제 보드의 역할 자격, 일별 용량, 가용 시간, 선후행, 마감, 변경 허용 범위.}{출력}{통과 후보 목록(중복 제거, 최대 50개). 탈락 후보는 어긴 제약을 기록한다.}};
\node[flow,below=of validate] (discuss) {\flowstep{05 | 팀원별 Proxy LLM: 후보 평가와 공개 토론}{입력}{통과 후보, 본인 위임카드(타인 카드는 볼 수 없음), 지금까지의 공개 발언.}{출력}{초기 평가 → 토론 대상 최대 10개 선정 → 공개 토론(질문·거절·조건) → 최종 평가.}};
\node[flow,below=of discuss] (consent) {\flowstep{06 | 변경 영향을 받는 팀원의 Proxy LLM: 후보별 동의}{입력}{후보 ID와 계획 hash로 지정한 정확한 계획, 본인 카드.}{출력}{후보마다 ACCEPT / REJECT / NEEDS\_INPUT. REJECT가 없는 후보 중 상위 1--3개를 PM 추천안으로 확정한다.}};
\node[flow,fill=amber,below=of consent] (pm) {\flowstep{07 | PM(사람): 추천안 중 하나를 직접 선택}{입력}{추천안별 변경 후 보드, 변경 내역, 공개 사유, 동의 상태.}{출력}{하나 선택 / 전부 반려 / 수정 요청. 추천안이 1개여도 자동 선택하지 않으며, PM 선택으로 팀원의 REJECT를 뒤집을 수 없다.}};
\node[flow,fill=amber,below=of pm] (human) {\flowstep{08 | NEEDS\_INPUT 팀원(사람): 본인 확인}{대상}{카드의 위임 범위를 넘는 변경, 또는 카드에 ``본인 확인''으로 지정된 변경을 받은 팀원만.}{출력}{모두 승인하면 09로 간다. 대상자가 없으면 이 단계를 건너뛴다.}};
\node[flow,fill=green,below=of human] (commit) {\flowstep{09 | Hive + 저장 코드: 최종 재검증 후 보드 반영}{검사}{snapshot 이후 보드·카드 변경 여부, 동의 대상과 선택 후보의 계획 hash 일치, 저장 후 보드의 제약 통과.}{출력}{모두 통과하면 한 트랜잭션으로 반영한다(\texttt{applied}).}};
\foreach \a/\b in {input/estimate,estimate/generate,generate/validate,validate/discuss,discuss/consent,consent/pm,pm/human,human/commit} \draw[arr] (\a)--(\b);
\node[side,right=6mm of validate] (none) {\sidetext{통과 후보가 0개}{탈락 사유를 PM에 전달하고 종료(\texttt{search\_exhausted}). 탐색 예산 안에서 못 찾았다는 뜻이며 해가 없다는 증명은 아니다. 마감 연기·인력 추가안은 미구현.}};
\draw[branch] (validate.east)--(none.west);
\node[side,right=6mm of discuss] (supp) {\sidetext{기존 후보로 못 푸는 조건이 나옴}{그 조건을 넣어 후보를 1회 추가 생성 → 04 검사 → 05 재평가. 기존 통과 후보는 그대로 둔다.}};
\draw[branch] (discuss.east)--(supp.west);
\node[side,right=6mm of consent] (nocons) {\sidetext{REJECT 없는 후보가 0개}{PM 검토로 넘기고 종료(\texttt{no\_consented\_candidate}).}};
\draw[branch] (consent.east)--(nocons.west);
\node[side,right=6mm of pm] (revise) {\sidetext{PM이 선택하지 않음}{수정 요청: PM 사유를 조건으로 추가해 03부터 다시 실행. 전부 반려: 보드를 바꾸지 않고 종료.}};
\draw[branch] (pm.east)--(revise.west);
\node[side,right=6mm of human] (deny) {\sidetext{팀원이 본인 확인을 거절}{그 후보를 제외한다. 남은 추천안이 있으면 07에서 PM이 다시 고르고, 없으면 PM 검토로 넘긴다.}};
\draw[branch] (human.east)--(deny.west);
\node[side,right=6mm of commit] (stop) {\sidetext{재검증 실패}{보드·카드가 바뀌었거나 저장 결과가 제약을 어기면 전부 롤백(부분 저장 없음)하고 PM에 사유를 전달한다.}};
\draw[branch] (commit.east)--(stop.west);
\end{tikzpicture}
\end{center}
{\footnotesize\textbf{그림 1.} 현행 V2 \texttt{candidate\_discussion} 경로. 01--09는 정상 경로, 오른쪽은 예외 경로다. 02·05·06은 개인 카드에 따른 LLM 판단, 04·09는 코드 검증으로, 두 판단은 서로를 대체하지 않는다.}

\newpage
\subsection{그림 1의 판단과 산출물 읽기}
\textbf{03의 산출물은 무엇인가?} 후보 하나는 담당자 명단이 아니라 새로운 일정 보드다. 원 Task의 범위를 어떤 산출물로 나누는지, 누가 언제 수행하는지, 어떤 조각이 어떤 입력을 기다리는지, 이후 milestone이 언제 끝나는지를 함께 담는다. 생성 AI는 의미 구조를 설계하고, 알고리즘은 담당·시간 조합을 확장한다. 이 단계에는 개인 비공개 위임카드를 전달하지 않는다.

\textbf{04가 통과했다고 합의인가?} 아니다. Hive 통과는 코드에 표현된 제약을 만족한다는 뜻이다. 예를 들어 Art가 일정을 수행할 수 있어도 자신의 위임 설정 때문에 새 인수를 거절할 수 있다. 그 판단은 05·06에서 Art의 Proxy가 한다. 반대로 LLM이 수락해도 Hive가 마감을 위반한다고 판단하면 적용할 수 없다.

\textbf{05·06에서 공정성은 어떻게 다루는가?} 각자의 기존 업무, 가용 시간, 맡게 될 실제 범위와 공개 조건을 근거로 후보를 비교한다. 부담 분산 수치는 참고 자료다. 현재 구현은 영향 당사자의 최종 평가에 거절이 없는 후보를 추리고, 참가자 점수 합계로 순서를 정한 뒤 후보별 동의를 별도로 받는다. 따라서 집계 규칙의 한계도 있으며, 단순 최소 분산을 공정성의 정답으로 보지 않는다.

\textbf{07·08은 왜 따로 있는가?} PM의 계획 선택권과 개인의 위임 권한이 다르기 때문이다. PM은 팀 차원의 후보를 최종 선택하고, 본인 확인은 해당 개인의 자동 위임 범위를 넘는 변경을 확인한다. NEEDS\_INPUT은 본인 확인 필요를 뜻해야 하며, 다른 사람의 승인 대기를 본인의 확인 사유로 혼동하는 Live 오류는 현재 남아 있는 제한이다.

\textbf{09에서 무엇을 다시 검사하는가?} 토론·사람 확인 중 보드나 카드가 바뀌지 않았는지, 선택한 후보와 동의 대상이 같은 버전인지, 실제 적용 결과도 마감과 제약을 지키는지를 검사한다. 다른 후보에 대한 옛 수락을 재사용하거나, 미합의 변경을 일부만 저장하지 않는다.

\textbf{T4로 읽는 예시.} Art의 Character Animation 증가로 최종 FX가 늦어졌다고 하자. 생성 단계는 Dev의 Integration을 임시 반영과 최종 교체로 나누고, FX 입력을 최종 교체에만 연결하는 후보를 만든다. Hive는 실제 Dev 가용 시간과 후속 QA 마감을 검사한다. Art·Dev·필요 시 QA는 자기 카드로 검토한다. Retest를 생략하는 후보라면 Final QA의 동일 항목 포함과 QA 본인 확인이 필요하다. Retest 유지 후보가 가능하면 그것도 유효 대안이다. PM이 선택하고 필요한 확인·최종 검증을 마쳐야 적용된다.

\textbf{모든 경로에 적용되는 종료 규칙.} 호출·시간 한도, 모델 장애, 후보 고갈, 동의 실패를 서로 다른 사유로 기록한다. 최대 토론 라운드는 최종 평가로 넘어가는 기준이며, 그 자체가 합의나 적용 성공을 뜻하지 않는다. 각 턴의 주체·수신자·후보·사유·Hive 검증과 보드 snapshot은 run 기록에서 확인할 수 있다.

\newpage
\section{시간 안에 가능한 새로운 milestone 설계}
후보는 단순히 ``누가 몇 일을 맡는가''가 아니다. 실제 산출물 범위, 담당자, 작업 순서, 중간 산출물, 의존관계, milestone 시점을 함께 담는 계획이다.

\subsection{생성 AI와 알고리즘의 역할 분리}
생성 AI에는 이슈, Task 설명·완료 조건·분리 경계, 공개 역할과 일정 정보를 전달한다. 개인 비공개 위임카드는 전달하지 않는다. 생성 결과는 구조화된 \texttt{Designs / Transformation}이며, 각 조각은 \texttt{key}, \texttt{what}, \texttt{effort}, \texttt{depends\_on}을 갖는다. 알고리즘은 이 구조를 실제 담당자와 시간 조합으로 확장한다. 명시된 Task 계약을 이용하는 결정적 템플릿도 활용한다.

\begin{tabularx}{\linewidth}{@{}p{28mm}Y@{}}\toprule
변환 & 계획에서 바뀌는 내용\\\midrule
분해·전체 양도 & 의미 있는 하위 산출물을 여러 명이 맡거나 Task 전체를 다른 사람에게 이전\\
일정 이동 & 허용된 미착수 업무의 시작 순서를 바꿔 후속 업무의 시간을 확보\\
협업 & 실제 지원 업무를 배정하고 주 작업 단축의 근거를 남김. Task에서 금지하면 차단\\
후속 이관 & 명시된 \texttt{defer\_option}에 따라 선택 범위를 Beta 등으로 이관. 소요량을 삭제하지 않음\\
선택 작업 생략 & optional 계약, 대체 검증 범위·최소 시간·당사자 확인을 충족할 때만 허용\\
단계 재구성 & 임시 반영·최종 교체처럼 명시된 하위 범위를 만들고 입력 의존성을 해당 단계에 연결\\\bottomrule
\end{tabularx}

분해 시 원 업무 기준 \texttt{effort}와 담당자별 실제 소요일 \texttt{days}를 구분한다. 수치만 나눈 조각은 의미 있는 업무 설계를 보장하지 않는다. 명시 계약은 코드로 검사하되, 범위의 의미상 적절성은 생성 근거와 당사자 토론에도 의존한다.

\subsection{Hive의 검증과 탐색 한계}
모든 후보는 격리된 보드에서 실행해 검사한다. 역할·허용 담당자, 일별 용량·가용 시간, 선행 산출물, Handoff·최종 마감, 업무량 보존, 허용 변경 범위를 검증한다. 사람의 승인이 필요한 0일 milestone과 자동 완료 장벽도 구별한다.

같은 실제 담당·일정·의존구조의 계획은 제목이나 임시 ID가 달라도 중복 제거한다. 공정성 점수는 해석과 토론의 참고이며, 생성 단계에서 최소 분산안을 자동 승자로 선택하지 않는다.

현재 기본 설정은 raw 조합 500개, 유효 후보 최대 50개, shortlist 최대 10개다. 후보 부족은 수학적 불가능성의 증명이 아니다. \texttt{search\_exhausted}는 해당 구조·탐색 예산 안에서 유효 후보를 확보하지 못했다는 뜻이다.

\newpage
\section{Discussion 확대: 개인 판단이 합의로 이어지는 과정}
\textbf{공통 입력.} 같은 후보 풀과 공개 Task·일정·이력. \textbf{개인 입력.} 호출 대상 본인의 카드와 역할 MD만 추가한다. 다음 발언자는 직전 발언을 포함한 갱신된 공개 대화를 받는다.
\begin{center}
\begin{tikzpicture}[node distance=6mm]
\node[flow] (assess) {\flowtext{A | 각 개인 LLM: 초기 평가}{후보를 최대 8개씩 읽고 본인의 업무·카드로 판단.}{후보별 의사, 점수, 공정성·카드 적합성 평가.}};
\node[flow,fill=green,below=of assess] (short) {\flowtext{B | 코드: 토론 후보 목록 구성}{개인별 상위 추천을 순회해 참가자 추천 슬롯 보장.}{shortlist를 고정하고 공개 토론 시작.}};
\node[flow,below=of short] (speak) {\flowtext{C | 이번 발언자 LLM: 공개 발언}{본인 카드 + shortlist + 질문·거절·조건의 전체 이력.}{수신자·후보 ID·행동·메시지·조건을 포함한 Speech.}};
\node[flow,fill=green,below=of speak] (route) {\flowtext{D | 코드: 발언 기록과 다음 호출}{후보·수신자 ID 검사. 공개 발언 저장, 새 수신자 초대.}{다음 사람에게 갱신된 이력 전달. 다음 라운드 순서 회전.}};
\node[flow,below=of route] (rank) {\flowtext{E | 각 개인 LLM: 토론 후 최종 평가}{새 조건과 상대 사유를 반영해 shortlist를 다시 비교.}{영향 당사자 REJECT가 있는 후보는 추천에서 제외.}};
\node[flow,below=of rank] (consent) {\flowtext{F | 영향 당사자 LLM: 후보별 Consent}{정확한 후보 ID·계획 hash로 본인의 변경 전체 확인.}{ACCEPT / REJECT / NEEDS\_INPUT과 공개 가능한 사유.}};
\node[flow,fill=amber,below=of consent] (pm) {\flowtext{G | 코드 → PM: 최종 후보 제출}{당사자 거절이 없는 후보를 최대 3개까지 추천.}{PM 선택 대기. 개인 LLM 점수가 자동 적용을 결정하지 않음.}};
\foreach \a/\b in {assess/short,short/speak,speak/route,route/rank,rank/consent,consent/pm} \draw[arr] (\a)--(\b);
\node[side,right=9mm of short] (speech) {\sidetext{실제로 무엇을 말하는가?}{QUESTION 질문 / ANSWER 응답 / PREFER 선호 / REJECT 거절 / CONDITION 조건 / PASS 추가 발언 없음.}};
\draw[branch] (speak.east)--(speech.west);
\node[side,right=9mm of route] (loop) {\sidetext{계속 토론하거나 후보 보충}{라운드가 남으면 C로 돌아간다. 새 후보 필요 조건은 보충 생성·Hive 검증 후 재평가한다. 동일 발언은 중복 표시를 억제하나 호출 자체는 이미 발생한 상태다.}};
\draw[branch] (route.east)--(loop.west);
\node[side,right=9mm of consent] (change) {\sidetext{새 계획에는 새 동의}{수정 후보의 일정·범위·영향 당사자를 다시 검사한다. 옛 계획의 수락은 새 계획의 동의가 아니다. 본인 확인은 PM 선택 뒤 별도다.}};
\draw[branch] (consent.east)--(change.west);
\end{tikzpicture}
\end{center}
{\footnotesize\textbf{그림 2.} A·C·E·F는 실제 개인 LLM 판단, B·D·G는 코드 처리다. 한 라운드는 여러 발언 호출로 이루어지며, 최종 평가와 후보별 동의는 별도 호출이다. 각 호출은 run 전체 예산을 공유한다.}

\textbf{이 과정의 한계.} 공개 발언이 실제 질문에 답했는지, 카드 해석이 일관적인지는 ID 검사만으로 보장되지 않는다. 같은 발언 반복을 건너뛰는 것과 의미 있는 합의를 형성하는 것도 다르다. 실제 Live 대화와 본인 확인 사유를 별도로 평가해야 한다.

\newpage
\section{PM 선택, 본인 확인, 안전한 적용}
정상적인 추천 완료는 \texttt{awaiting\_pm\_selection}이다. 후보가 하나여도 PM이 명시적으로 선택해야 한다. PM은 전체 반려와 수정 요청도 할 수 있다. 후보가 0개인 실패 전달 \texttt{waiting\_pm}과 구분한다.

PM 선택에는 후보 ID, plan hash와 보드 snapshot을 결합한다. 권한 밖 위임 또는 명시 확인 대상은 \texttt{NEEDS\_INPUT}으로 본인 확인을 기다린다. PM 선택이 개인의 거절을 수락으로 바꾸지 않는다. 확인 대기 중 보드나 카드가 바뀌면 기존 선택을 그대로 적용하지 않고 다시 검사한다. 적용은 트랜잭션 안에서 수행하고 실제 저장 결과도 Hive로 확인한다.

\begin{tabularx}{\linewidth}{@{}p{46mm}Y@{}}\toprule
상태·종료 사유 & 해석\\\midrule
\texttt{running} & 후보 생성·평가·토론·동의 진행 중\\
\texttt{awaiting\_pm\_selection} & 정상 추천 완료, PM 최종 선택 대기\\
\texttt{waiting\_confirm} & 선택안에 필요한 당사자 사람 확인 대기\\
\texttt{applied} & 필요한 동의와 최신 검증 후 실제 보드 반영\\
\texttt{search\_exhausted} & 제한된 탐색에서 후보를 확보하지 못함\\
\texttt{no\_consented\_candidate} & 최종 당사자 비거절 후보가 남지 않음\\
모델·호출·시간 한도 종료 & 모델 오류와 후보 부재를 구분하고 PM에게 요약 전달\\\bottomrule
\end{tabularx}

기본 discussion은 2라운드이며 설정 범위는 1--15다. run 전체에는 기본 120회 실제 호출 예산과 600초 한도를 둔다. 전송 재시도도 예산에 포함한다. 생성 응답 형식 오류는 구체적인 필드 오류를 주고 한 번 복구한다. 평가 후보 ID가 잘못되면 기대 ID·누락·잘못된 ID를 명시해 한 번 다시 요청하며 재실패 시 중단한다. 반복 오류를 무제한 재시도하지 않는다.

\section{화면, API와 기록}
제공된 \texttt{front/Plan Bee Front.html}과 어댑터가 FastAPI에 연결된다. 화면에서 PM / Dev / Art / GD / QA 역할을 바꿔 체험한다. 예시 이슈 버튼은 입력창만 채우며, 보내기를 눌러야 실행한다. 역할 전환 자체는 LLM을 호출하지 않는다.

\begin{tabularx}{\linewidth}{@{}p{69mm}Y@{}}\toprule
API & 목적\\\midrule
\texttt{POST /candidate-runs} & 업무와 줄글 이슈로 run 시작\\
\texttt{POST /runs/\{id\}/advance} & 저장된 상태에서 다음 단계 진행\\
\texttt{POST /runs/\{id\}/pm-selection} & PM 선택·반려·수정 요청\\
\texttt{POST /confirmations/\{id\}} & 해당 본인의 확인 또는 거절\\
\texttt{GET /state}, \texttt{GET /board} & 현재 보드와 제약 결과\\\bottomrule
\end{tabularx}

실행 기록은 \texttt{runs/v2/live/}와 \texttt{runs/v2/mock/}에 날짜·시간·고유값 파일명으로 분리한다. 초기·턴별·최종 milestone, 라운드, 발화자·수신자, 후보와 사유, Hive 검사, 동의 상태·종료 사유를 추적한다. 공개 기록과 비공개 카드·평가 저장을 구분한다. 화면의 역할 헤더는 로컬 체험용 제어이며 실제 로그인 인증 체계를 대신하지 않는다.

\newpage
\section{에이전트별 MD 프롬프트의 역할}
\texttt{prompts/} 아래 MD는 각 LLM 호출의 지시문이다. 코드를 바꾸지 않고 에이전트의 판단 기준을 조정하는 곳이며, run 시작 시 hash로 고정되어 실행 중 수정이 섞이지 않는다.

\textbf{한 번의 개인 Proxy 호출 구성.} system 메시지는 \texttt{discussion/system.md}(공통 규칙) + \texttt{discussion/actor.md}(본인 이름·위임카드) + \texttt{agents/}(역할)\texttt{.md}(본인 역할 지침)이다. user 메시지는 호출별 MD(assess, speak, rank, candidate\_consent 등)와 입력 데이터다. 다른 사람의 역할 MD와 카드는 들어가지 않으며, 생성 AI 호출에는 역할 MD와 카드가 모두 없다.

\begin{tabularx}{\linewidth}{@{}p{36mm}Y@{}}\toprule
역할 MD & 판단 기준\\\midrule
\texttt{agents/dev.md}\newline Dev (구현) & 원래 소요량 보존 확인(단순 이동에 소요 증가 금지). Task의 분리 경계(로직 전담 / 에셋 연결) 준수. 기능 우선순위는 GD에게 공개 질문. 임시 반영과 최종 교체 구분. QA 시간 단축으로 충돌을 덮지 않음.\\
\texttt{agents/art.md}\newline Art (에셋) & 가능한 에셋 연결은 검토. 본인 확인이 필요한 신규 인수는 NEEDS\_INPUT이며 금지(REJECT)와 구분. 지연 시 기존·새 총량, 최종 전달 시각, 중간 산출물을 구분. 카드 원문 미인용.\\
\texttt{agents/gd.md}\newline GD (기획·범위) & 카드의 기능 우선순위와 범위 권한을 따름. 필수 Core 보호, 선택 기능 이관이 완료 기준을 유지하는지 검토. 근거만 공개하고, 자기 동의가 PM 선택을 대신하지 않음.\\
\texttt{agents/qa.md}\newline QA (검증) & 선행 완료 전 검증이나 Final QA 단축에 동의하지 않음. 선택 중간 Retest와 필수 최종 검증 구분. 허용된 Retest 생략도 본인 확인이 필요하면 NEEDS\_INPUT. 실제 검증 시작·종료를 설명.\\
\texttt{agents/reestimate.md}\newline 담당자 재견적(02) & 일정 이동인지 소요량 변화인지(duration\_changed) 먼저 판단. 소요가 바뀌면 새 총 견적, 근거, 이슈 원문 인용을 답함.\\\bottomrule
\end{tabularx}

\vspace{4mm}
\begin{tabularx}{\linewidth}{@{}p{52mm}Y@{}}\toprule
호출별 MD (그림 1 단계) & 지시 내용\\\midrule
\texttt{candidate\_generation/system}, \texttt{transform} (03) & 카드 없이 이슈와 Task 계약으로 서로 다른 계획 구조를 최대 12개 설계. 승자를 고르지 않고 목적·완료 기준·마감 유지.\\
\texttt{candidate\_generation/repair}, \texttt{format\_repair} (03) & 모든 후보가 Hive에서 탈락하면 사유를 근거로 구조를 1회 재설계. 형식 오류는 필드별 오류를 받아 1회 재제출.\\
\texttt{discussion/assess}, \texttt{rank} (05) & 모든 후보에 ACCEPT / REJECT / NEEDS\_INPUT과 종합·공정성·카드 적합성 점수. rank는 토론 내용을 반영.\\
\texttt{discussion/speak} (05) & QUESTION, ANSWER, PREFER, REJECT, CONDITION, PASS 중 하나로 공개 발언. 앞 발언에 실제로 응답.\\
\texttt{discussion/candidate\_consent} (06) & 정확한 후보 ID·plan hash의 전체 변경에 대한 당사자 동의.\\
\texttt{discussion/*\_repair} & 후보 ID 누락·오류나 스키마 오류를 구체적 피드백과 함께 1회 재요청.\\\bottomrule
\end{tabularx}

\newpage
\section{사람 순위와의 별도 비교 평가}
운영 PM 선택과 블라인드 사람 순위 평가는 분리한다. 고정된 동일 후보 집합에서 공정성과 본인 카드 적합성을 별도 축으로 비교한다. 사람과 LLM에게 제공되는 카드·공개 정보 조건을 맞추고, 사람에게 에이전트 선호를 먼저 노출하지 않는다.

후보 $i$의 사람 순위를 $r_h(i)$, 에이전트 순위를 $r_a(i)$라 할 때 rank-softmax로 분포를 만든다. 동률은 평균 순위로 처리한다.
\[
p_h(i)=\frac{\exp(-(r_h(i)-1)/\tau)}{\sum_j\exp(-(r_h(j)-1)/\tau)},\qquad
q_a(i)=\frac{\exp(-(r_a(i)-1)/\tau)}{\sum_j\exp(-(r_a(j)-1)/\tau)}.
\]
\[
\operatorname{CE}(p_h,q_a)=-\sum_i p_h(i)\log q_a(i).
\]
기본 $\tau=1$이며 낮은 CE가 사람 분포에 더 가깝다. 사람 entropy와 KL, 후보 수, 관측 coverage, 균등 baseline $\log K$, Kendall $\tau_b$ 등을 함께 해석한다. CE는 순위 전체의 순열 likelihood가 아니다. 미평가 후보를 임의로 최하위 처리하지 않으며, 후보가 하나면 순위 비교를 수행하지 않는다. 실제 사람 라벨 없이 점수를 만들어 보고하지 않는다.

\section{현재 연결된 다섯 시나리오}
\begin{tabularx}{\linewidth}{@{}p{13mm}p{30mm}Y@{}}\toprule
구분 & 핵심 이슈 & 검토할 계획·승인\\\midrule
T1 & 일정 이동 & 6시간 Polish를 다음 월요일로 이동, 소요량 불변·Dev 중심 검토\\
T2 & 업무 재분배 & UI 로직은 Dev, 에셋 연결은 Art. Art 신규 인수에 본인 확인\\
T3 & 범위 이관 & 선택 기능과 연결된 잔여 FX를 Beta로 이관, 필수 Core·QA 유지\\
T4 & 복합 조정 & 임시 Integration / 최종 교체와 입력 의존성 재구성. 선택 Retest 생략 시 QA 확인\\
T5 & 현재 제약 충돌 & 필수 수정 뒤 QA 시간을 확보할 후보가 없으면 Dev·QA 근거와 PM 검토 전달\\\bottomrule
\end{tabularx}
원문은 \texttt{seed/m4\_t*/source.md}, 초기 보드·Task·카드는 \texttt{team.yaml}, 이슈는 \texttt{feed.yaml}에 둔다. 원문의 시간 불일치를 정규화한 부분은 시나리오 문서에 공개한다. T4는 Retest를 유지하면서 기한을 충족하는 유효 대안도 있으므로 생략을 강제하지 않는다.

\newpage
\section{T1--T5 기대 결과와 실제 Live 실행 비교}
같은 판정기(\texttt{scripts/check\_m4\_results.py})로 프롬프트·코드 버전별 Live 실행(gpt-4o-mini, 시나리오당 1회)을 채점했다. 사람의 PM 선택과 본인 확인은 시뮬레이션하지 않으므로, ``도달''은 적용 성공이 아니라 기대한 상태와 기대 후보에 도달했다는 뜻이다.

\begin{tabularx}{\linewidth}{@{}p{20mm}p{25mm}Y@{}}\toprule
시나리오 & 기대 상태 & 기대 후보\\\midrule
T1 일정 이동 & PM 선택 대기 & Polish만 다음 월요일로 이동, 소요 0.75일 유지, Bug Fix 불변. 참가자 Dev 한 명.\\
T2 재분배 & PM 선택 대기 & 로직 1일(Dev) / 에셋 연결 1일(Art) 분담, 후속 일정 유지. Art 본인 확인.\\
T3 범위 이관 & PM 선택 대기 & Advanced Effects 4일 + 잔여 FX 0.5일 Beta 이관, Core·Final QA 유지, GD 영향 당사자.\\
T4 복합 조정 & PM 선택 대기 & Integration 임시/최종 2단계(합 1일), 최종 교체는 FX 이후. Retest 생략안은 QA 확인.\\
T5 해결 불가 & PM 검토 & 유효 후보 0개(\texttt{search\_exhausted}), Dev·QA 검토 의견 전달, 마감 연장 없음.\\\bottomrule
\end{tabularx}

\vspace{4mm}
\begin{tabularx}{\linewidth}{@{}p{20mm}YYYY@{}}\toprule
 & v006 & v007 & v008 & v009 (최신)\\\midrule
T1 & 도달 (5회, 13초) & 도달 (6회, 16초) & 도달 (5회, 12초) & 도달 (5회, 11초)\\
T2 & 도달 (13회, 34초) & 도달 (14회, 32초) & 도달 (13회, 29초) & 도달 (13회, 26초)\\
T3 & 도달 (18회, 58초) & \textbf{생성 형식 오류} & \textbf{토론 참조 오류} & 도달 (17회, 44초)\\
T4 & \textbf{평가 ID 오류} & 중단 & 도달 (21회, 57초) & 도달 (18회, 57초)\\
T5 & 도달 (7회, 34초) & 미실행 & 도달 (5회, 22초) & 도달 (5회, 19초)\\\bottomrule
\end{tabularx}

{\footnotesize 괄호는 LLM 호출 수와 소요 시간. v007 실험은 T4 토론 중 사용자 요청으로 중단되어 T5를 실행하지 않았다. 원본: \texttt{past\_versions/agent\_prompts/}(버전)\texttt{/experiments/*\_live/}.}

\subsection{차이 분석}
\textbf{실패는 모두 LLM 응답 형식에서 발생했다.} 일정·제약 판단(Hive)은 모든 버전에서 기대와 같았고, 미도달은 모델 응답이 스키마나 참조 규칙을 어긴 지점에서 run이 종료된 경우다. 버전마다 그 지점을 하나씩 보완했다.
\begin{itemize}
\item v006 → v007: T4 평가가 존재하지 않는 후보 ID를 참조. 정확한 ID 목록으로 1회 재요청하도록 보완.
\item v007 → v008: T3 생성에서 프롬프트의 ``이관은 parts=[]''를 모델이 template 최상위 parts로 오해해 응답 전체가 탈락. 빈 parts를 원 업무 전체로 정규화하고, 형식이 틀린 template 하나만 제외하도록 변경.
\item v008 → v009: T3 생성은 통과(기대 후보 3개)했으나 공개 발언이 목록에 없는 후보·수신자를 가리켜 종료. 수신자 이름을 id로 정규화하고 후보 ID를 1회 재요청.
\end{itemize}

\textbf{기대와 일치한 부분.} T5는 매번 후보 0개로 PM 검토에 전달되었고 마감을 연장하지 않았다. T4(v009)는 기대에 맞는 후보 4개 중 Retest 유지안과 QA 확인이 필요한 생략안이 함께 나와, 생략을 정답으로 강제하지 않는 시나리오 의도와 맞았다. T1은 Dev 한 명만 참여해 호출 5회로 끝났다.

\textbf{해석의 한계.} 시나리오당 1회 실행이라 성공률이 아니다. v009에서도 T2 평가 재요청, T3 template 제외가 실제로 발생했으며, 같은 버전에서도 응답에 따라 다른 지점에서 오류가 날 수 있다. PM 선택·본인 확인·실제 적용 성공은 별도로 평가해야 한다.

\newpage
\section{Ablation: 단순 알고리즘 / 토론 제외 / 전체}
\textbf{목적.} 생성 AI와 개인 카드 토론이 각각 무엇을 더하는지, 그리고 결과가 시나리오에 미리 구조화된 정보에 얼마나 의존하는지 확인한다. 같은 T1--T5, 같은 판정기, 같은 seed로 실행했다(\texttt{scripts/ablation\_m4.py}, 조건·시나리오당 Live 1회, gpt-4o-mini).

\begin{tabularx}{\linewidth}{@{}p{27mm}Y@{}}\toprule
조건 & 구성\\\midrule
단순 알고리즘 & LLM 0회. 이슈를 읽지 못하므로 새 소요일만 구조화 값으로 제공(유리한 가정). 코드 기본 구조만 탐색하고 부담 분산이 가장 작은 Hive 통과안을 자동 적용.\\
토론 제외 & 재견적·생성 AI·Hive까지. 평가·토론·동의 없이 분산이 작은 상위 3개를 PM에 제출.\\
전체 (Plan Bee) & 생성 AI + Hive + 개인 카드 평가·토론·후보별 동의 + PM 선택.\\\bottomrule
\end{tabularx}

\textbf{두 가지 시드.} \emph{구조화 시드}는 원래 시나리오다. \emph{이슈 힌트 제거}는 이슈의 해결 방향을 미리 적어 둔 필드(T1 \texttt{requested\_start}·\texttt{single\_actor}, T2 \texttt{required\_parts}·\texttt{split\_boundary})를 지워, 해결 방향을 이슈 원문에서만 알 수 있게 했다. T3 이관 권한과 T4 단계 분해·검증 대체 계약은 이슈와 무관한 사전 권한이라 남겼다.

\subsection{결과: 프롬프트 개선 전(v009)과 후(v010)}
\begin{tabularx}{\linewidth}{@{}p{10mm}YYY|YYY@{}}\toprule
 & \multicolumn{3}{c|}{구조화 시드 (v010)} & \multicolumn{3}{c}{이슈 힌트 제거: v009 → v010}\\
 & 알고리즘 & 토론 제외 & 전체 & 알고리즘 & 토론 제외 & 전체\\\midrule
T1 & 도달 & 도달 & 도달 & 미도달 & 미도달 → \textbf{도달} & 미도달 → \textbf{도달}\\
T2 & 도달\textsuperscript{a} & 도달\textsuperscript{b} & 도달\textsuperscript{c} & 후보 0\textsuperscript{g} & 후보 0\textsuperscript{g} & 후보 0\textsuperscript{g}\\
T3 & 도달 & 도달 & 도달 & 도달 & 도달 → 도달 & 도달 → 도달\\
T4 & 도달 & 도달\textsuperscript{d} & 도달 & 도달 & 도달\textsuperscript{d} → 도달\textsuperscript{d} & 동의 0\textsuperscript{e} → \textbf{도달}\\
T5 & 부분\textsuperscript{f} & 부분\textsuperscript{f} & 도달 & 부분\textsuperscript{f} & 부분 → 부분 & 도달 → 도달\\\midrule
LLM & 0회 & 2--3회 & 5--17회 & 0회 & 2--3회 & 3--17회\\\bottomrule
\end{tabularx}

{\footnotesize ``도달'' = 판정기 기준 기대 계획이 결과(알고리즘: 자동 적용 1개, 그 외: PM 추천안)에 포함. 구조화 시드는 v009와 v010 결과가 같다.
a: Art 본인 확인이 필요한 분담안을 묻지 않고 자동 적용. b: 같은 안을 동의 정보 없이 PM에 제출. c: Art Proxy가 NEEDS\_INPUT으로 본인 확인 요구.
d: QA 확인이 필요한 Retest 생략안을 동의 없이 추천안에 포함. e: Dev·QA Proxy가 Hive 통과 후보를 사실과 다른 ``일정 겹침''으로 거절.
f: 후보 0개는 맞지만 Dev·QA 검토 의견 없이 종료. g: 남은 권한 안에서 유효한 계획이 없어 PM 검토로 전달(아래 2번). 원본: \texttt{eval/ablation/}, \texttt{past\_versions/agent\_prompts/*v010*}.}

\subsection{원인 분석}
\textbf{1. 구조화 시드에서는 알고리즘만으로도 기대 계획에 닿았다.} Task에 해결 방향과 권한 계약이 구조화되어 있으면 Hive 탐색만으로 정답 구조를 만든다. 이 설정에서 LLM의 차별점은 정답 찾기가 아니라 동의 처리(각주 a--d)와 검토 의견(f)이었다. 비용은 LLM 호출 2--3회 대비 5--17회다.

\textbf{2. T2의 필드는 힌트이자 권한이었다.} \texttt{ui\_integration}의 담당 가능자는 Dev뿐이고, 다른 역할(Art)이 조각을 맡으려면 Task 계약(\texttt{required\_parts})이 필요하다(``역할을 넘는 분담은 AI 주장이 아니라 명시 계약으로만''). 이 필드를 지우면 Dev 단독 2일안만 남아 마감을 어기므로, 후보 0개로 PM에 넘기는 것이 의도된 동작이다. 처음 실험에서 이 필드를 순수한 이슈 힌트로 분류한 것은 잘못이었다.

\textbf{3. T1 실패는 생성 지시의 구조 문제였다.} v009 생성 응답을 직접 확인한 결과, 모델은 (i) 프롬프트에 나열된 defer·restructure·omit를 템플릿마다 하나씩 채우는 ``메뉴''로 읽었고, (ii) 원 업무를 두고 시작만 옮기는 기본안을 내지 않았으며, (iii) 6h를 effort 6으로 쓰고, (iv) 일부 구조에서 필수 \texttt{basis}를 빠뜨렸다. 권한 계약이 없는 연산은 Hive가 전부 탈락시키므로 이동 후보가 남지 않았다.

\textbf{4. T4 오거절은 토론 지시의 빈틈이었다.} Proxy가 Hive가 이미 검증한 일정을 다시 의심해 거절했고, 다른 참여자의 사유를 복사했다.

\subsection{개선 (v010)}
\begin{itemize}
\item \texttt{candidate\_generation/transform.md}: ``이슈 유형 판별 → 원 업무 유지 기본안(필수) → 이슈가 말한 대안'' 절차로 재작성. 시작 시각 계산(현재 s + 이동량), 일 단위 환산 예시, T1형 정답 형식 예시를 추가. plan\_actions는 해당 권한 키가 있을 때만 쓰는 예외로 명시.
\item \texttt{candidate\_generation/repair.md}: 반려 원인별 복구 방향(병렬 분해, 다른 미착수 업무 이동, 권한 없는 연산 제거, 분량 보존)을 구체화.
\item \texttt{discussion/system.md}(모든 개인 Proxy 공통): Hive 통과 후보는 일정 겹침·마감을 거절 사유로 쓰지 않고, 우려는 구체 시각을 인용한 조건·질문으로. 거절은 본인 카드 근거로만, 타인 사유 복사 금지.
\item 생성 후처리(\texttt{generator.py}): 프롬프트로 금지해도 모델이 분해를 권한 없는 \texttt{restructure}로 표현하는 습관이 남아, 권한 계약이 없는 plan\_actions 항목만 제거하고 나머지 구조를 살린다. 남은 구조도 Hive 전체 검증을 그대로 받으므로 권한을 우회하지 않는다.
\end{itemize}

\textbf{효과와 한계.} 이슈 힌트 없이도 T1에서 요청한 이동 후보가 생성되어 기대 결과에 도달했고, 구조화 시드의 결과는 그대로 유지되었다(후보 경로 테스트 29개 통과). T4 전체 조건의 회복은 토론 지시 개선의 효과일 수 있으나 1회 실행이라 실행 간 변동과 구분할 수 없다. T2는 권한 설계상 프롬프트로 해결할 대상이 아니며, 실제 운용에서는 이슈에서 역할 간 분담 계약을 추출해 PM·당사자가 승인하는 단계가 필요하다. 모든 칸은 1회 실행이며 성공률로 해석하지 않는다.

\section{구현 근거와 현재 제한}
주요 구현은 \texttt{app/candidates/engine.py}의 상태 머신, \texttt{generator.py}의 후보 확장·중복 제거, \texttt{plans.py}의 preview·적용, \texttt{app/hive.py}의 일정·제약 검사다. 개인 호출 어댑터는 \texttt{app/llm.py}, 역할별 MD는 \texttt{prompts/agents/}, 생성·토론 MD는 \texttt{prompts/candidate\_generation/}, \texttt{prompts/discussion/}에 있다.

v010 기준 전체 회귀 테스트 205개가 통과했다(그중 V2 후보 경로 테스트 29개). Live에서는 v009로 T1--T5 모두 기대 상태에 도달했고(10절), v010은 구조화 시드에서 같은 결과를 유지하면서 이슈 힌트 없는 T1도 해결했지만(11절), 시나리오당 1회 실행이며 형식 오류 복구가 실제로 발생했다. 불필요한 본인 확인 해석 같은 카드 해석 문제도 별도 평가가 필요하다. 따라서 모든 Live 시나리오가 안정적으로 성공한다고 주장하지 않는다.

후보 의미의 완전성, 제한 탐색의 해 공간 커버리지, LLM의 카드 해석 일관성은 여전히 한계다. 과거 프롬프트와 실제 실험은 \texttt{past\_versions/agent\_prompts/}에 보존한다. PM 대기 도달, 사람 승인 완료, 실제 적용 성공, 사람 선호 일치도는 서로 다른 지표로 보고해야 한다.
\end{document}
